% Generated by scripts/assemble.py; corrections belong in source/corrections.json.
\examhead{183}{西北工业大学 2026 年数学分析试卷}{172}
% p173-b20; source PDF page 173
\par\Needspace{5\baselineskip}
\blockanchor{p173-b20}
\majorquestion{一}（20分）用极限的严格数学定义证明：\par
% p173-b21; source PDF page 173
\blockanchor{p173-b21}
(1) $\lim_{x\to +\infty}\frac{x^2 + x - 1}{x^2} = 1.$\par
% p173-b22; source PDF page 173
\blockanchor{p173-b22}
(2) $\lim_{n\to \infty}\sqrt[n]{a} = 1$ ，其中 $a > 1$\par
% p173-b23; source PDF page 173
\par\Needspace{5\baselineskip}
\blockanchor{p173-b23}
\majorquestion{二}（20分）解答如下问题：\par
% p173-b24; source PDF page 173
\blockanchor{p173-b24}
(1) 求极限: $\lim_{x \to +\infty} \left( \sqrt{x^4 + 1} - (x^2 - 1)e^{\frac{1}{x}} \right)$.\par
% p173-b25; source PDF page 173
\blockanchor{p173-b25}
(2) 求幂级数 $\sum_{n=1}^{\infty} \frac{x^n}{n(n+1)2^n}$ 的收敛域与和函数.\par
% p173-b26; source PDF page 173
\par\Needspace{5\baselineskip}
\blockanchor{p173-b26}
\majorquestion{三}(15 分) 使用闭区间套定理证明聚点定理: 实数轴上任一有界无限点集 $S$ 必有一个聚点.\par
% p174-b00; source PDF page 174
\par\Needspace{5\baselineskip}
\blockanchor{p174-b00}
\majorquestion{四}(15 分) 证明: 若 $f(x)$ 在 $[a, b]$ 上黎曼可积, 则 $\lim_{p \to +\infty} \int_{a}^{b} f(x) \cos px \, \mathrm{d}x = 0$.\par
% p174-b01; source PDF page 174
\par\Needspace{5\baselineskip}
\blockanchor{p174-b01}
\majorquestion{五}（15分）利用 $\int_0^{+\infty}e^{-x^2}\mathrm{d}x = \frac{\sqrt{\pi}}{2}$ ，计算 $\int_0^{+\infty}\frac{e^{-a^2x^2} - e^{-b^2x^2}}{x^2}\mathrm{d}x,$ 其中 $b > a > 0.$\par
% p174-b02; source PDF page 174
\par\Needspace{5\baselineskip}
\par\Needspace{6\baselineskip}
\blockanchor{p174-b02}
\majorquestion{六}(15 分) 函数 $z = z(x, y)$ 二阶偏导连续, 试以 $u, v$ 作为新的自变量变换方程\par
\nopagebreak[4]
% p174-b03; source PDF page 174
\blockanchor{p174-b03}
\displayformula{x ^ {2} \frac {\partial^ {2} z}{\partial x ^ {2}} - y ^ {2} \frac {\partial^ {2} z}{\partial y ^ {2}} = 0.}
% p174-b04; source PDF page 174
\blockanchor{p174-b04}
其中 u = xy,  $v = \frac{x}{y}$ .\par
% p174-b05; source PDF page 174
\par\Needspace{5\baselineskip}
\blockanchor{p174-b05}
\majorquestion{七}(15 分) 证明: $f(x) = \sqrt{x}$ 在 $[0, +\infty)$ 上一致连续.\par
% p174-b06; source PDF page 174
\par\Needspace{5\baselineskip}
\par\Needspace{6\baselineskip}
\blockanchor{p174-b06}
\majorquestion{八}（15分）应用高斯公式计算三重积分\par
\nopagebreak[4]
% p174-b07; source PDF page 174
\blockanchor{p174-b07}
\displayformula{\iiint_ {V} (x y + y z + z x) \mathrm{d} x \mathrm{d} y \mathrm{d} z.}
% p174-b08; source PDF page 174
\blockanchor{p174-b08}
其中 $V$ 是由 $x \geq 0, y \geq 0, 0 \leq z \leq 1$ 和 $x^2 + y^2 \leq 1$ 所确定的空间区域.\par
% p174-b09; source PDF page 174
\par\Needspace{5\baselineskip}
\blockanchor{p174-b09}
\majorquestion{九}(15 分) 函数 $f(x)$ 在 $(- \infty, + \infty)$ 二阶可导, $f(x)$ 是 $(- \infty, + \infty)$ 上的有界函数, 证明: 存在 $\xi \in (-\infty, + \infty)$, 满足 $f''(\xi) = 0$.\par
% p174-b10; source PDF page 174
\par\Needspace{5\baselineskip}
\par\Needspace{6\baselineskip}
\blockanchor{p174-b10}
\majorquestion{十}(5 分) 函数 $f(x)$ 在 $(- \infty, +\infty)$ 上各阶可导, 且存在 $M > 0$, 使得\par
\nopagebreak[4]
% p174-b11; source PDF page 174
\blockanchor{p174-b11}
\displayformula{\left| f ^ {(k)} (x) \right| \leq M, \forall x \in (- \infty , + \infty), k = 1, 2, \dots .}
% p174-b12; source PDF page 174
\blockanchor{p174-b12}
如果在一无限有界集 E 上,  $f(x) \equiv 0$ , 证明: 在  $(-∞, +∞)$  上,  $f(x) \equiv 0$ .\par
